\section{JavaScript}

\begin{frame}[fragile]{JavaScript in Odoo 17}
	\begin{itemize}
		\item Odoo frontend code is organized as ES modules.
		\item Files usually start with:
\begin{lstlisting}[style=js]
/** @odoo-module **/
\end{lstlisting}
		\item Modules import from Odoo packages such as:
		\begin{itemize}
			\item \texttt{@odoo/owl}
			\item \texttt{@web/core/...}
			\item \texttt{@web/views/...}
		\end{itemize}
	\end{itemize}
\end{frame}

\begin{frame}[fragile]{Module Import / Export}
\begin{lstlisting}[style=js]
/** @odoo-module **/
import { Component, useState } from "@odoo/owl";
import { registry } from "@web/core/registry";

export class StudentCounter extends Component {
	static template = "school_manager.StudentCounter";
	setup() {
		this.state = useState({ value: 0 });
	}
	increment() {
		this.state.value++;
	}
}
\end{lstlisting}
\end{frame}

\begin{frame}[fragile]{Registering Frontend Pieces}
\begin{lstlisting}[style=js]
registry.category("actions").add("school_student_dashboard", StudentDashboard);
registry.category("fields").add("student_badge", studentBadgeField);
\end{lstlisting}
	\begin{itemize}
		\item Registries are how Odoo discovers actions, fields, views, services, and more.
		\item Custom JS usually extends an existing registry category.
	\end{itemize}
\end{frame}

\begin{frame}[fragile]{Debugging Tips}
	\begin{itemize}
		\item Use browser devtools extensively.
		\item Check Assets for load/order issues after upgrades.
		\item Prefer small modules over giant JS files.
		\item Keep naming consistent with module technical names.
	\end{itemize}
\end{frame}

\begin{frame}{JavaScript Best Practices}
	\begin{itemize}
		\item Always include \texttt{/** @odoo-module **/}.
		\item Prefer official \texttt{@web} / \texttt{@odoo/owl} imports.
		\item Avoid editing core Odoo JS; extend via your module assets.
		\item Write intentional UI logic, not copies of backend business rules.
	\end{itemize}
\end{frame}
